% Capítulo 1: Introducción
% Propósito: Presentar la introducción, el planteamiento del proyecto, los objetivos, los antecedentes, la justificación, el alcance, las limitaciones y la viabilidad del proyecto.
% Contenido: tomado del documento "TRABAJO DE CULMINACIÓN DE GRADO - AVANCES".

\section{Introducción}

En Centroamérica, las comunidades costeras dependen operativamente de centros de acopio pesqueros que articulan la compra de mariscos en muelle con el mercado local. Sin embargo, la persistencia de métodos analógicos, como bitácoras en cuadernos y anotaciones en pizarras, limita la eficiencia operativa y genera desincronizaciones severas entre el inventario físico y las ventas. Como señalan \textcite{ahumada2020sistema}, la falta de un registro transaccional en tiempo real proporciona inconsistencias de datos, descuadres financieros y una pérdida progresiva de trazabilidad en la cadena de suministro.

El Acopio Arrecife es uno de los establecimientos formalmente acreditados por el INPESCA en San Juan del Sur, Rivas, y enfrenta esta problemática en su día a día. Aunque adquiere producto marino diariamente y lo conserva en termos y congeladores, gestiona sus preventas e inventarios mediante mensajes de WhatsApp, llamadas y registros manuales que se actualizan únicamente hasta la entrega del producto. Esta desconexión provoca el riesgo constante de sobreventa (overbooking) de especies previamente apartadas y dificulta la obtención de métricas comerciales precisas sobre los ingresos y las especies de mayor demanda.

Para resolver esta limitación, el presente proyecto plantea el desarrollo de una plataforma web propia y a medida bajo el patrón Clean Architecture. De acuerdo con \textcite{villasmil2021}, la adopción de interfaces web transaccionales en MiPymes elimina el margen de error en el conteo de stock y asegura el bloqueo dinámico de productos en tiempo real. La plataforma permitirá registrar los lotes desde su ingreso, congelar el stock ante preventas, ofrecer modalidades de cobro flexible y generar reportes semanales automáticos, todo bajo la administración autónoma del personal del propio acopio.

\subsection{Planteamiento del Proyecto}

En las comunidades costeras de Centroamérica, el sector pesquero artesanal e intermedio depende operacionalmente de centros de acopio que actúan como articuladores comerciales entre los pescadores y el mercado gastronómico local. En el municipio de San Juan del Sur, Rivas, el Acopio «Arrecife» adquiere diariamente especies marinas en el muelle pesquero, trasladándolas a su local para su conservación en congeladores (freezers) y termos con hielo.

A pesar de estar formalmente acreditado por el Instituto Nicaragüense de la Pesca y Acuicultura (INPESCA), el Acopio Arrecife ejecuta la comercialización mediante herramientas analógicas, tales como la exposición de oferta en pizarras acrílicas externas, el registro financiero en cuadernos y la gestión de ventas o preventas mediante llamadas telefónicas o mensajería instantánea. De acuerdo con \textcite{ahumada2020sistema}, la dependencia de registros manuales en papel induce desincronizaciones en el control de stock, duplicidad de datos y una marcada falta de trazabilidad en las existencias reales.

Esta ineficiencia se agudiza cuando los clientes realizan apartados o preventas por teléfono o mensajería. Al no contar con un sistema automatizado que reserve el stock en tiempo real, las preventas no se descuentan inmediatamente de la pizarra ni del cuaderno, sino hasta la entrega física del producto. Esta situación desencadena el riesgo constante de sobreventa (overbooking) al ofrecer productos previamente apartados, genera descuadres entre las libras reales en freezers, termos y las registradas, y provoca la ausencia de métricas comerciales al dificultar la generación de reportes consolidados sobre ingresos y especies más vendidas sin realizar cálculos manuales propensos a errores.

Como señalan \textcite{mayorga2024} y \textcite{villasmil2021}, la implementación de aplicaciones web centralizadas bajo bases de datos transaccionales elimina la brecha de información operativa, optimiza la rotación del inventario y automatiza los reportes financieros en pequeñas y medianas empresas.

\subsection{Objetivos}

\subsubsection{Objetivo General}

Desarrollar una plataforma web a medida que automatice el control de inventario en tiempo real, el registro de preventas y la generación de reportes comerciales para el Acopio «Arrecife» en el municipio de San Juan del Sur, durante el período 2027-2028.

\subsubsection{Objetivos Específicos}

\begin{enumerate}
  \item Analizar los procesos operativos actuales de recepción de producto marino y flujo de preventas en el Acopio «Arrecife» para la especificación formal de requisitos del software.
  \item Definir los requisitos de software funsionales y no funsionales.
  \item Diseñar la arquitectura lógica del sistema, el modelo de base de datos y las interfaces de usuario (UI/UX) para el control de inventario en tiempo real, preventas y reportes comerciales.
  \item Desarrollar los módulos funcionales de la plataforma web (API REST y Frontend) que permitan el ingreso de lotes, el apartado inmediato de stock por preventa, la confirmación de entregas mediante cobro flexible y la emisión automática de resúmenes de ventas.
  \item Validar el funcionamiento de la plataforma web en el Acopio «Arrecife» mediante pruebas transaccionales de inventario y evaluaciones de usabilidad con el personal del establecimiento para asegurar su operación autónoma.
\end{enumerate}

\subsection{Antecedentes}

A nivel internacional, diversos estudios científicos han abordado la digitalización de procesos comerciales e inventarios mediante tecnologías web. \textcite{villasmil2021}, en su investigación titulada \emph{Sistema web para la gestión e integración de procesos de inventario en MiPymes} y publicada en la Revista de Ciencias Tecnológicas, desarrollaron una plataforma distribuida para el control de inventario en tiempo real. En dicho estudio demostraron que la migración de un control analógico a una interfaz web transaccional redujo en un 82\,\% el margen de error en el conteo de stock y eliminó eficazmente las pérdidas por sobreventa.

Desde una perspectiva metodológica, \textcite{ahumada2020sistema}, en una investigación publicada en la Revista ESPACIOS, evaluaron el impacto de las bases de datos relacionales en la sustitución de bitácoras manuales en cuadernos. Los autores concluyeron que la persistencia de datos relacional asegura la consistencia de las transacciones mediante las propiedades de la atomicidad, consistencia, aislamiento y durabilidad (ACID), previniendo así cualquier inconsistencia o duplicidad de datos cuando se efectúan múltiples pedidos de forma simultánea.

A nivel regional, el estudio realizado por la organización \textcite{goalglobal2021} en la costa del Caribe centroamericano (a través del programa MiPesca) implementó un entorno web de trazabilidad y gestión para la red de organizaciones de pescadores artesanales y sus centros de acopio (RED-PESCAHN). La investigación concluyó que los sistemas ligeros con interfaces intuitivas permitieron a los encargados de acopio registrar entradas y salidas de lotes sin requerir asistencia técnica constante, logrando formalizar el canal de comercialización y preventas con restaurantes y mercados regionales.

A nivel nacional, los diagnósticos oficiales sobre la cadena de valor pesquera del Instituto Nicaragüense de la Pesca y Acuicultura, \textcite{inpesca2023}, señalan que la mayoría de los centros de acopio artesanales, tanto en el litoral Pacífico como en la Costa Caribe (RACCN y RACCS), operan mediante bitácoras de papel, recibos manuscritos y pizarras. Mientras que en el Caribe el registro en papel se vincula al aislamiento geográfico de las faenas y al control por cooperativas, en el Pacífico predomina la anotación manual diaria en locales de playa. En ambos contextos, esta dependencia analógica genera descuadres en la consolidación del peso por lote y retrasos en la actualización de las existencias para los compradores, evidenciando la necesidad de implementar soluciones web transaccionales adaptadas al contexto local.

\subsection{Justificación}

El Acopio «Arrecife», ubicado en el municipio de San Juan del Sur y formalmente acreditado ante el Instituto Nicaragüense de la Pesca y Acuicultura \parencite{inpesca2023}, realiza diariamente la adquisición de productos marinos en el muelle pesquero para su posterior conservación en recipientes de termos y freezer. No obstante, su gestión comercial se ejecuta mediante métodos empíricos y analógicos: la exposición estática de precios en pizarras acrílicas, el registro contable en cuadernos y la atención de pedidos mediante llamadas telefónicas o mensajería instantánea (WhatsApp). Como señalan \textcite{ahumada2020sistema}, la dependencia de registros manuscritos carentes de persistencia transaccional centralizada genera una marcada ineficiencia operativa, inconsistencia de datos y un desgaste diario en la administración de pequeños centros de acopio.

Esta limitación se intensifica cuando los clientes habituales, tales como restaurantes, hoteles y hogares, solicitan apartados o preventas. Al no contar con un mecanismo de sincronización inmediata, el pedido no se descuenta del cuaderno ni se retira de la pizarra sino hasta la entrega física del producto, desencadenando la duplicidad de reservas y el riesgo constante de sobreventa (overbooking). Al respecto, \textcite{villasmil2021} afirman que el desfase entre las consultas de los compradores y la disponibilidad física real incrementa la insatisfacción del cliente y ocasiona pérdidas económicas por mermas transaccionales. Asimismo, la consolidación manual de ventas en libretas imposibilita la obtención oportuna de métricas comerciales sobre el volumen de ingresos semanales y las especies de mayor rotación, afectando la toma de decisiones estratégicas \parencite{laudon2020}.

Para superar estas barreras operativas, se fundamenta la creación de una plataforma web propia y a medida para el Acopio «Arrecife», la cual automatizará el flujo del producto desde su recepción física en el local. El sistema permitirá registrar el ingreso de lotes por termos o freezer, bloquear dinámicamente el stock en tiempo real ante solicitudes de preventa y ofrecer modalidades de cobro flexible, ya sea efectivo o verificación de transferencias bancarias directas. De acuerdo con \textcite{mayorga2024}, la implementación de arquitecturas web centralizadas en MiPymes comerciales elimina el margen de error en el conteo de existencias y automatiza la generación de reportes financieros analíticos sin recurrir a cálculos manuales propensos a fallas.

Un atributo indispensable de esta solución tecnológica radica en su diseño orientado a la usabilidad y la autogestión, lo que garantizará que el propio personal del acopio opere y administre la herramienta de forma autónoma sin depender de alguna persona de afuera especializada \parencite{davis1989,goalglobal2021}. En el ámbito regional, estudios aplicados en la cadena de valor pesquera artesanal demuestran que la sustitución de bitácoras de papel por entornos web sencillos optimiza la rotación de inventarios, elimina las pérdidas por sobreventa y eleva la competitividad comercial del negocio \parencite{goalglobal2021,pressman2021}. En consecuencia, el proyecto dotará al Acopio «Arrecife» de una herramienta informática robusta, ágil y adaptada a su contexto cotidiano, impulsando su modernización operativa en San Juan del Sur.

\subsection{Alcance}

El desarrollo de la plataforma web para el Acopio "Arrecife" contempla una arquitectura modular, esto significa que su infraestructura no se va a construir como un bloque único e inamovible, sino a partir de componentes independientes, estandarizados y combinables, orientada a la automatización del flujo comercial del negocio. Los módulos que tendremos en cuenta son los siguientes: 

. Módulo de Inventario en Tiempo Real: permitirá la captura y registro exacto de los lotes pesqueros al momento de su recepción, especificando la especie, el peso en libras y el contenedor asignado (termos con hielo o congeladores). De acuerdo con las guías de trazabilidad en centros de acopio (FEDEPESCA, 2020), la digitalización inmediata de entradas por contenedor es la base indispensable para garantizar la consistencia del stock físico frente a las consultas de los compradores. 

. Módulo de Gestión de Preventas: encargado de registrar al instante los apartados solicitados por los clientes mediante llamadas o mensajería instantánea. Como señalan Villasmil et al. (2021), la implementación de reservas temporales con bloqueo de stock en bases de datos transaccionales resulta indispensable para eliminar el riesgo de sobreventa (overbooking) y descontar automáticamente las existencias visibles para el público. 

. Módulo de Confirmación de Venta y Pago: gestionará la actualización del estado de las órdenes a Entregado/Pagado, admitiendo modalidades de cobro flexible mediante efectivo o verificación de transferencias y depósitos bancarios directos. 

. Módulo de Reportes Comerciales: generación automática de resúmenes semanales y mensuales de ventas e ingresos netos por especie, proporcionando métricas clave para la toma de decisiones gerenciales. \textcite{Laudon, Kenneth C and Laudon.2012}.

. Módulo de Mantenimiento Autónomo: dotado de una interfaz administrativa intuitiva, diseñada bajo principios de usabilidad para permitir que el propio personal del Acopio "Arrecife" opere y administre la plataforma sin depender de soporte técnico continuo.


\subsection{Limitaciones}

El desarrollo del proyecto presenta una serie de delimitaciones metodológicas y tecnológicas indispensables para acotar la ejecución del software durante el período 2027-2028. Las limitaciones estimadas son las siguientes:

. Ámbito de Aplicación: la investigación se limita exclusivamente al modelado y automatización de los procesos operativos internos del Acopio "Arrecife" en San Juan del Sur, por lo que las reglas de negocio y los flujos transaccionales no pretenden ser extrapolables a otros centros de acopio con dinámicas distintas (Ahumada et al., 2020).

. Módulo de Pagos y Facturación: el sistema no procesará cobros automáticos mediante pasarelas con tarjetas de crédito o débito internacionales en esta primera fase. En su lugar, el modelo de cobro flexible adoptado limitará la validación de pagos por transferencia o depósito bancario a la verificación manual de comprobantes por parte del encargado del acopio antes de marcar la orden como "Pagada", un enfoque recomendado para MiPymes pesqueras artesanales que buscan reducir costos de integración y comisiones financieras en etapas iniciales de digitalización (GOAL Global, 2021).

. Infraestructura tecnológica y Conectividad: la plataforma requiere de una conexión activa a Internet, ya sea en el local (wifi) o en el dispositivo del encargado, para asegurar la sincronización transaccional del inventario y la actualización inmediata del stock en tiempo real. Como señalan Villasmil et al. (2021), la persistencia e integridad de los datos en aplicaciones web transaccionales dependen de la conectividad continua con la base de datos central para prevenir la concurrencia de pedidos desactualizados.


\subsection{Viabilidad del Proyecto}

\subsubsection{Viabilidad Económica}

La factibilidad económica del proyecto se evaluó mediante un análisis de costo beneficio proyectado a un horizonte temporal de tres años, metodología recomendada para la incorporación de tecnologías de información en pequeñas y medianas empresas \parencite{kendall2014,maldonado2021}. El presupuesto de inversión inicial ($I_{0}$) asciende a \$1,400.00 USD, monto que contempla los honorarios profesionales por el diseño, codificación y pruebas de la plataforma bajo el patrón Clean Architecture, la adquisición del dominio web, la configuración inicial del servidor virtual privado (VPS) y la capacitación del personal del negocio. Por su parte, los costos operativos anuales se estiman en \$120.00 USD destinados al mantenimiento del hosting en la nube y licencias operativas. Al contrastar estos egresos con la eficiencia económica generada por el sistema, estimada en \$840.00 USD anuales debido a la eliminación de mermas por sobreventa, la prevención de descuadres transaccionales y la optimización de los tiempos de atención, se obtiene un flujo neto de efectivo de \$720.00 USD al año.

Para determinar la viabilidad de la inversión, se aplicaron los indicadores financieros clásicos de decisión considerando una tasa de descuento ($k$) del 12\,\% \parencite{lledo2007}. El cálculo del Valor Actual Neto (VAN) arrojó un resultado positivo de +\$329.32 USD, lo que demuestra que la implementación del software genera valor económico agregado por encima del costo de oportunidad del capital \parencite{burneoValarezo2016}. Asimismo, la Tasa Interna de Retorno (TIR) alcanzó un 25.1\,\%, superando significativamente la tasa de descuento de referencia y ratificando la rentabilidad del proyecto. Finalmente, el Período de Recuperación de la Inversión (PRI) se determinó en 1.94 años (aproximadamente un año y once meses), confirmando que la empresa recuperará la totalidad del capital invertido antes del segundo año de operación. Por consiguiente, desde la perspectiva financiera, el proyecto resulta plenamente viable y ejecutable.

\subsubsection{Viabilidad Técnica}

Desde la perspectiva técnica, la viabilidad de la propuesta se sustenta en la adopción de una arquitectura de software robusta, desacoplada y basada en estándares de código abierto. La plataforma se construirá utilizando el patrón Clean Architecture con una API REST en el backend desarrollada en .NET e interfaz de usuario en Angular, utilizando PostgreSQL como motor de base de datos relacional \parencite{elmasri2020,villasmil2021}. Esta infraestructura no requiere el desembolso de licencias propietarias y garantiza un alto grado de mantenibilidad \parencite{sommerville2019}. Asimismo, la solución aprovechará la conectividad a Internet de banda ancha y redes móviles presente en el municipio de San Juan del Sur, permitiendo que el encargado del acopio acceda a la plataforma transaccional a través de navegadores web en dispositivos móviles (smartphones) o computadoras convencionales sin necesidad de adquirir hardware especializado adicional.

\subsubsection{Viabilidad Operativa}

En el ámbito operativo, la factibilidad se fundamenta en el Modelo de Aceptación Tecnológica (TAM), el cual establece que la percepción de utilidad y la facilidad de uso determinan la adopción efectiva de un sistema \parencite{davis1989}. La plataforma web ha sido diseñada para alinearse con los procesos cotidianos del Acopio «Arrecife», tales como la recepción en muelle, la conservación en termos o congeladores y el cobro flexible en efectivo o transferencia, sustituyendo el registro manual en cuadernos y pizarras por una interfaz táctil simplificada que previene errores humanos \parencite{ahumada2020sistema}. La automatización del apartado de preventas por WhatsApp o llamada bloquea el stock de forma inmediata, eliminando la duplicidad de pedidos e incentivando la aceptación por parte del personal \parencite{villasmil2021}. Además, la sencillez de los módulos garantiza que la administración y actualización del catálogo sean gestionadas de forma autónoma por los propios miembros del acopio, asegurando la sostenibilidad de la herramienta en el tiempo.

\subsubsection{Viabilidad Legal}

La implementación de la plataforma web se adecua estrictamente al marco jurídico nicaragüense aplicable al comercio, la protección de datos e inocuidad en el sector pesquero. En primer lugar, la digitalización de los volúmenes de captura, especies y almacenamiento responde a los requerimientos de trazabilidad exigidos por la Ley No. 489 (Ley de Pesca y Acuicultura de Nicaragua) y las normativas dictadas por el Instituto Nicaragüense de la Pesca y Acuicultura (INPESCA), facilitando el cumplimiento legal del establecimiento en San Juan del Sur \parencite{republicaNicaragua2004}. Asimismo, las funcionalidades de registro de ventas e inventario contribuyen al mantenimiento de estándares de higiene e inocuidad regulados por el Ministerio de Salud (MINSA).

En lo que concierne al manejo de información y transacciones electrónicas, el sistema garantiza el cumplimiento de la Ley No. 787 (Ley de Protección de Datos Personales) mediante el cifrado de información de clientes, proveedores y comprobantes bancarios a través de protocolos de seguridad HTTPS/SSL \parencite{republicaNicaragua2012}. A su vez, la verificación de comprobantes de transferencias y la emisión de comprobantes digitales encuentran respaldo normativo en la Ley No. 729 (Ley de Firma Electrónica y Comercio Electrónico). Finalmente, la plataforma se desarrollará como una solución informática a medida para el acopio conforme a la Ley No. 312 (Ley de Derechos de Autor y Derechos Conexos), utilizando componentes y marcos de trabajo bajo licencias libres (MIT y BSD), lo que libera a la empresa de cualquier conflicto de propiedad intelectual o contingencia legal por uso de software no autorizado.
